% !TEX root = ../../ctfp-print.tex

\lettrine[lhang=0.17]{我}{们已经看到两种基本}组合类型的机制：使用乘积和余积。事实上，日常编程中的许多数据结构都可以仅通过这两种机制构建。这一事实具有重要的实际意义。数据结构的许多性质都是可组合的。例如，如果你知道如何比较基本类型的值是否相等，并且知道如何将这些比较推广到乘积类型和余积类型，你就可以自动推导复合类型的等式运算符。在Haskell中，你可以为一大类复合类型自动推导等式、比较、字符串转换等功能。

让我们更仔细地考察编程语言中的乘积类型与和类型。

\section{乘积类型}

编程语言中两个类型的乘积的标准实现是一个二元组(pair)。在Haskell中，二元组是一个原始类型构造器；而在C++中，它是一个相对复杂的模板，定义在标准库中。

\begin{figure}[H]
  \centering
  \includegraphics[width=0.35\textwidth]{images/pair.jpg}
\end{figure}

\noindent
二元组并非严格满足交换律：一个\code{(Int, Bool)}类型的二元组不能替代\code{(Bool, Int)}类型的二元组，尽管它们携带相同的信息。不过它们在同构意义下满足交换律——这种同构由\code{swap}函数给出（该函数是自身的逆函数）：

\src{snippet01}
你可以把这两个二元组看作只是以不同格式存储相同数据。这就像大端序(big endian)与小端序(little endian)的区别。

你可以通过嵌套二元组的方式组合任意数量的类型形成乘积，但存在更简单的方法：嵌套二元组等价于多元组(tuple)。这是由于不同的嵌套方式之间存在同构关系。如果你想按顺序组合三种类型\code{a}、\code{b}和\code{c}形成乘积，可以通过两种方式实现：

\src{snippet02}

或
\src{snippet03}
这些类型是不同的——你不能将一种类型传递给期望另一种类型的函数——但它们的元素之间存在一一对应关系。存在一个映射两者的函数：

\src{snippet04}
并且这个函数是可逆的：

\src{snippet05}
因此这是一个同构关系。这些只是对相同数据的不同打包方式。

你可以将乘积类型的创建解释为类型上的二元运算。从这个角度看，上述同构非常类似于我们在幺半群(monoid)中见过的结合律：
$$(a * b) * c = a * (b * c)$$
只不过在幺半群的情况下，这两种组合方式是严格相等的，而在这里它们只在"同构意义下相等"。

如果我们接受同构的概念而不坚持严格的相等性，可以进一步证明单位类型\code{()}在乘积运算中的作用如同乘法中的1一样是单位元。确实，将某个类型\code{a}的值与单位类型配对不会增加任何信息。类型：

\src{snippet06}
与\code{a}是同构的。以下是具体的同构关系：

\src{snippet07}

\src{snippet08}
这些观察结论可以通过以下形式化表述：$\Set$（集合范畴）是一个\newterm{幺半范畴}(monoidal category)。这个范畴同时也具有类似单子(monoid)的结构，在此意义上你可以对对象进行"乘法"操作（即取笛卡尔积）。我将在后续讨论幺半范畴并给出完整定义。

在Haskell中存在一种更通用的方式来定义乘积类型——特别是当我们即将看到的与和类型结合使用时。这种方式使用带有多个参数的命名构造器(named constructors)。例如，二元组可以替代性地定义为：

\src{snippet09}
其中，\code{Pair a b}是参数化于另外两个类型\code{a}和\code{b}的类型名称；\code{P}是数据构造器的名称。通过向类型构造器\code{Pair}传递两个类型来定义二元组类型；通过向构造器\code{P}传递两个适当类型的值来构造二元组值。例如，定义一个\code{String}与\code{Bool}的组合值\code{stmt}：

\src{snippet10}
第一行是类型声明，使用了类型构造器\code{Pair}，其中\code{String}和\code{Bool}分别替换了泛型定义中的\code{a}和\code{b}。第二行通过向数据构造器\code{P}传入具体字符串和布尔值来定义实际值。类型构造器用于构造类型，数据构造器用于构造值。

由于Haskell中类型构造器和数据构造器的命名空间相互独立，你经常会看到同一名称同时用于两者的情况，例如：

\src{snippet11}
如果你眯着眼睛看，甚至可以把内置的二元组类型视为这类声明的一种变体，其中名称\code{Pair}被二元运算符\code{(,)}取代。实际上你可以像使用其他命名构造器一样使用\code{(,)}，并通过前缀记法创建二元组：

\src{snippet12}
类似地，可以用\code{(,,)}创建三元组，依此类推。

除了使用泛型二元组或多参构造器外，你还可以定义具有特定名称的乘积类型，例如：

\src{snippet13}
这本质上就是\code{String}与\code{Bool}的乘积类型，但赋予了独立的名称和构造器。这种声明风格的优势在于，你可以定义多个具有相同内容但不同语义和功能的类型，这些类型之间不可互换。

使用元组和多参构造器进行编程可能会变得混乱且容易出错——需要时刻追踪每个组件代表什么。通常更优的做法是为组件命名。在Haskell中，具有命名字段的乘积类型称为\newterm{记录}(record)，在C语言中则称为\code{struct}。

\section{记录}

让我们来看一个简单的例子。我们希望通过组合两个字符串（名称和符号）以及一个整数（原子序数）来描述化学元素，将它们组合成一个数据结构。我们可以使用元组\code{(String, String, Int)}并记住每个分量的含义。我们会通过模式匹配提取分量，例如下面这个检查元素符号是否为其名称前缀的函数（如\textbf{He}是\textbf{Helium}的前缀）：

\src{snippet14}
这段代码容易出错，且难以阅读和维护。更好的方法是定义一个记录：

\src{snippet15}
这两种表示方式是同构的，这可以通过这两个互为逆函数的转换函数来证明：

\src{snippet16}

\src{snippet17}
注意记录字段的名称同时也作为访问这些字段的函数。例如，\code{atomicNumber e}会从\code{e}中提取\code{atomicNumber}字段。我们使用\code{atomicNumber}作为具有如下类型的函数：

\src{snippet18}
通过\code{Element}的记录语法，我们的函数\code{startsWithSymbol}变得更具可读性：

\src{snippet19}
我们甚至可以使用Haskell技巧，通过将函数\code{isPrefixOf}用反引号包围转化为中缀运算符，使代码几乎像自然语句一样：

\src{snippet20}
在这种情况下可以省略括号，因为中缀运算符的优先级低于函数调用。

\section{和类型}

正如集合范畴中的乘积会导出乘积类型一样，余积(coproduct)则会导出和类型(sum type)。Haskell中和类型的规范实现是：

\src{snippet21}
同序对类似，\code{Either}具有交换性(至多差一个同构)，可以嵌套且嵌套顺序无关(至多差一个同构)。因此我们可以定义三元组对应的和类型：

\src{snippet22}
依此类推。

事实上，$\Set$范畴关于余积也构成一个(对称)幺半群范畴。其中二元运算由不交并实现，单位元由始对象担任。在类型层面，我们用\code{Either}作为幺半群运算符，用空类型\code{Void}作为单位元。可以将\code{Either}视为加法，\code{Void}视为零元素。确实，向和类型添加\code{Void}不会改变其内容。例如：

\src{snippet23}
就与\code{a}同构。这是因为无法构造该类型的\code{Right}版本——不存在\code{Void}类型的值。\code{Either a Void}的唯一实例通过\code{Left}构造器创建，仅封装了\code{a}类型的值。因此从符号上看，$a + 0 = a$。

和类型在Haskell中非常常见，但其C++对应物联合体(union)或变体(variant)却较少使用。原因有以下几点：

首先，最简单的和类型就是枚举类型，在C++中用\code{enum}实现。Haskell和类型：

\src{snippet24}
对应的C++代码为：

\begin{snip}{cpp}
enum { Red, Green, Blue };
\end{snip}
更简单的和类型：

\src{snippet25}
即C++的原始类型\code{bool}。

用于编码值存在性的简单和类型在C++中常通过特殊技巧和"不可能"值(如空字符串、负数、空指针等)实现。这种可选性若显式表达，在Haskell中使用\code{Maybe}类型：

\src{snippet26}
\code{Maybe}类型是两个类型的和。若将两个构造器分离为独立类型：
第一个构造器：

\src{snippet27}
这是一个名为\code{Nothing}的单值枚举。换句话说，它等价于单位类型\code{()}。第二个部分：

\src{snippet28}
仅封装了\code{a}类型。我们本可以将\code{Maybe}编码为：

\src{snippet29}
更复杂的和类型在C++中常用指针伪造。指针要么为空，要么指向特定类型值。例如Haskell列表类型可定义为递归和类型：

\src{snippet30}
在C++中可通过空指针技巧实现空列表：

\begin{snip}{cpp}
template<class A>
class List { 
    Node<A> * _head;
public:
    List() : _head(nullptr) {} // 空列表
    List(A a, List<A> l)       // 构造
      : _head(new Node<A>(a, l))
    {}
};
\end{snip}
注意Haskell的两个构造器\code{Nil}和\code{Cons}被转换为两个重载的\code{List}构造函数，参数分别对应空参数(用于\code{Nil})和值与列表的组合(用于\code{Cons})。\code{List}类不需要标签区分和类型的两个分量，而是通过\code{\_head}的特殊值\code{nullptr}编码\code{Nil}。

不过主要区别在于Haskell数据结构的不可变性。一旦使用某个构造器创建对象，该对象将永久记住所用构造器及其参数。因此\code{Maybe}对象若以\code{Just "energy"}创建，永远不会变成\code{Nothing}。同样，空列表将永远为空，三个元素的列表也将始终包含相同元素。

正是这种不可变性使得构造过程可逆。给定对象总能分解回其构造时使用的部件。这种解构通过模式匹配完成，复用构造器作为模式。构造器参数(若有)将被替换为变量(或其他模式)。

\code{List}数据类型有两个构造器，因此任意\code{List}的解构需要对应这两个构造器的模式。一个匹配空列表\code{Nil}，另一个匹配\code{Cons}构造的列表。例如，以下是\code{List}上的简单函数定义：

\src{snippet31}
\code{maybeTail}定义的第一部分使用\code{Nil}构造器作为模式返回\code{Nothing}。第二部分使用\code{Cons}构造器作为模式。第一个构造器参数被替换为通配符(因为我们不关心它)，\code{Cons}的第二个参数绑定到变量\code{t}(即使严格来说它们永不变化，我仍将这些称为变量：一旦绑定到表达式，变量永不改变)。返回值为\code{Just t}。现在，根据你的\code{List}创建方式，它将匹配其中一个子句。若使用\code{Cons}创建，则传递给它的两个参数将被提取(第一个被丢弃)。

更复杂的和类型在C++中通过多态类层次实现。具有共同祖先的一系类可以理解为一个变体类型，其中虚表(vtable)充当隐藏标签。在Haskell中通过构造器模式匹配和调用特化代码实现的功能，在C++中通过基于虚表指针派发虚拟函数调用来完成。

你很少看到C++中使用\code{union}作为和类型，因为联合体对可容纳成员有严重限制。甚至不能将\code{std::string}放入联合体，因为它具有拷贝构造函数。

\section{类型的代数}

单独来看，乘积类型和求和类型可以用来定义多种有用的数据结构，但真正的力量来源于两者的结合。我们再次调用了组合的力量。

让我们总结目前已知的内容。我们发现了两种底层的类型系统结构：求和类型具有单位元\code{Void}，乘积类型具有单位元单位类型\code{()}。我们可以将其类比为加法与乘法。在这个类比中，\code{Void}对应零，单位类型\code{()}对应一。

让我们看看这个类比能延伸多远。例如，零乘以任何数是否得到零？换句话说，一个包含\code{Void}分量的乘积类型是否与\code{Void}同构？例如，能否创建一个\code{Int}和\code{Void}的有序对？

创建有序对需要两个值。虽然很容易找到整数值，但不存在\code{Void}类型的值。因此，对于任意类型\code{a}，类型\code{(a, Void)}都是不可实例化的——没有值——因此等价于\code{Void}。换句话说，$a \times 0 = 0$。

将加法与乘法联系起来的另一个性质是分配律：
$$a \times (b + c) = a \times b + a \times c$$
这对乘积类型和求和类型是否成立？是的——如同往常一样，直到同构为止。左边对应类型：

\src{snippet32}
右边对应类型：

\src{snippet33}
这是单向转换的函数：

\src{snippet34}
这是反向转换的函数：

\src{snippet35}
\code{case of}语句用于函数内部的模式匹配。每个模式后跟箭头和匹配时要计算的表达式。例如，当用值：

\src{snippet36}
调用\code{prodToSum}时，\code{case e of}中的\code{e}等于\code{Left "Hi!"}。它会匹配模式\code{Left y}，将\code{"Hi!"}代入\code{y}。由于\code{x}已匹配为\code{2}，\code{case of}子句及整个函数的结果将是\code{Left (2, "Hi!")}，符合预期。

我不会证明这两个函数互为逆函数，但仔细思考的话，它们必然如此！它们只是简单地重新打包了两种数据结构的内容。本质上是相同的数据，只是格式不同。

数学家给这两种交织的幺半群起了个名字：称为\newterm{半环}(semiring)。它不是完整的\newterm{环}(ring)，因为我们无法定义类型的减法。这就是为什么半环有时被称为\newterm{rig}(pun on "ring without an 'n'")。抛开这点不谈，我们可以通过将自然数（构成rig的集合）的陈述翻译到类型论中获得很多启示。以下是几个有趣的对应关系：

\begin{longtable}[]{@{}ll@{}}
  \toprule
  自然数      & 类型\tabularnewline
  \midrule
  \endhead
  $0$          & \code{Void}\tabularnewline
  $1$          & \code{()}\tabularnewline
  $a + b$      & \code{Either a b = Left a | Right b}\tabularnewline
  $a \times b$ & \code{(a, b)} 或 \code{Pair a b = Pair a b}\tabularnewline
  $2 = 1 + 1$  & \code{data Bool = True | False}\tabularnewline
  $1 + a$      & \code{data Maybe = Nothing | Just a}\tabularnewline
  \bottomrule
\end{longtable}

\noindent
列表类型非常有趣，因为它被定义为某个方程的解。我们定义的类型会出现在方程两侧：

\src{snippet37}
如果我们进行常规替换并将\code{List a}记作\code{x}，就会得到方程：

\begin{Verbatim}
  x = 1 + a * x
\end{Verbatim}
我们无法用传统代数方法求解，因为不能对类型做减法或除法。但我们可以通过不断用右侧的\code{(1 + a*x)}替换\code{x}，并利用分配律得到以下序列：

\begin{Verbatim}
  x = 1 + a*x
  x = 1 + a*(1 + a*x) = 1 + a + a*a*x
  x = 1 + a + a*a*(1 + a*x) = 1 + a + a*a + a*a*a*x
  ...
  x = 1 + a + a*a + a*a*a + a*a*a*a...
\end{Verbatim}
最终得到无限项的乘积和（元组），可以解释为：列表要么是空的\code{1}；要么是单元素\code{a}；要么是二元组\code{a*a}；三元组\code{a*a*a}；等等。这正是列表的本质——一串连续的\code{a}！

列表还有更多特性，当我们学习函子和不动点后会再次回到它们和其他递归数据结构。

用符号变量解方程——这就是代数！这也正是这些类型名称的由来：代数数据类型(algebraic data types)。

最后，我应该提到类型代数的一个重要诠释。注意到两个类型\code{a}和\code{b}的乘积必须同时包含类型\code{a}和类型\code{b}的值，这意味着两个类型都必须可实例化。而两个类型的求和只需要包含类型\code{a}或类型\code{b}的值即可，只要其中一个可实例化就足够了。逻辑上的"与"(and)和"或"(or)同样构成半环，也可以映射到类型论：

\begin{longtable}[]{@{}ll@{}}
  \toprule
  逻辑                & 类型\tabularnewline
  \midrule
  \endhead
  $\mathit{false}$     & \code{Void}\tabularnewline
  $\mathit{true}$      & \code{()}\tabularnewline
  $a \mathbin{||} b$   & \code{Either a b = Left a | Right b}\tabularnewline
  $a \mathbin{\&\&} b$ & \code{(a, b)}\tabularnewline
  \bottomrule
\end{longtable}

\noindent
这个类比更加深刻，构成了逻辑与类型论之间Curry-Howard同构的基础。当我们讨论函数类型时会再次回到这个主题。

\section{挑战}

\begin{enumerate}
  \tightlist
  \item
        展示 \code{Maybe a} 和
        \code{Either () a} 之间的同构(isomorphism)。
  \item
        以下是用Haskell定义的一个和类型(sum type)：

        \begin{snip}{haskell}
data Shape = Circle Float 
           | Rect Float Float
\end{snip}
        当我们要定义像 \code{area} 这样作用于 \code{Shape} 的函数时，通常通过对两个构造器进行模式匹配(pattern matching)来实现：

        \begin{snip}{haskell}
area :: Shape -> Float
area (Circle r) = pi * r * r
area (Rect d h) = d * h
\end{snip}
        在C++或Java中将 \code{Shape} 实现为接口，并创建两个类：\code{Circle} 和 \code{Rect}。将 \code{area} 实现为虚函数(virtual function)。
  \item
        延续上例：我们可以轻松添加一个新的函数 \code{circ} 来计算 \code{Shape} 的周长(circumference)，且无需修改 \code{Shape} 的原始定义：

        \begin{snip}{haskell}
circ :: Shape -> Float
circ (Circle r) = 2.0 * pi * r
circ (Rect d h) = 2.0 * (d + h)
\end{snip}
        将 \code{circ} 添加到你的C++或Java实现中。你修改了原始代码的哪些部分？
  \item
        继续扩展：向 \code{Shape} 添加新形状 \code{Square} 并进行必要的更新。对比Haskell与C++/Java实现时，你分别修改了哪些代码？（即使你不是Haskell程序员，这种修改方式也应该是显而易见的。）
  \item
        证明对类型而言 $a + a = 2 \times a$ 在同构意义下成立（up to isomorphism）。请记住根据我们的映射表，$2$ 对应的是 \code{Bool} 类型。
\end{enumerate}